\section{Convergence and surjectivity of the coordinates}
\label{sec:uniformization}

We now pass from the polynomial models to a global map.  The local
errors must remain small under backward polynomial iteration, and the
subsequential forward degree bound must then rule out a proper image.
We give both steps with their domains and parameter choices.

\subsection{Parameter order and local error estimates}

Decrease the common chart exponent $\sigma$ of
Proposition~\ref{prop:shrinking-charts}, if necessary, so that
$0<\sigma<1$.  Choose
\begin{equation}
 0<\gamma<\gamma'<\min\{\sigma/2,1\},
 \qquad 0<10\nu<\min\{c\gamma,\sigma,\gamma,1\},
 \label{eq:unif:decay-parameters}
\end{equation}
with $\gamma$ in the range of its small-ball transition estimate;
decrease its constant $c$ to at most one.  Fix the jet order so large
that
\begin{equation}
 (d+1)\nu>30.
 \label{eq:unif:jet-order}
\end{equation}
Only after this choice fix a sufficiently small $h>0$, as specified
below.  Then choose $N,M,b,K_0,\lambda$ as in
Section~\ref{sec:dynamics}, and finally choose a sufficiently late
regular macro start.  The estimates in
Proposition~\ref{prop:polynomial-models} hold for this choice.

On a fixed smaller raw transition ball, the jet conjugacy and Cauchy
estimates give
\begin{equation}
 \abs{H_{k+1}F_k(z)-G_kH_k(z)}
 \le \exp\bigl(A_d\Delta s_k+o(s_k)\bigr)\abs z^{d+1}.
 \label{eq:unif:remainder}
\end{equation}
Here $A_d$ is independent of $h$.  Indeed the degrees of the
polynomials in this expression are fixed, the $F_k$ are uniformly
bounded on a larger raw ball, and the coefficient estimates of
Proposition~\ref{prop:polynomial-models} bound the two compositions
there.  They agree through degree $d$, so the Cauchy remainder has
the displayed factor.  This argument uses $H_{k+1}F_k-G_kH_k$
directly; it requires no uniform analytic domain for $H_k^{-1}$.

We shall use two sequences of raw chart coordinates:
\begin{enumerate}[label=(\roman*)]
 \item $z_k=f_{s_k}(x)$, uniformly for $x$ in any fixed compact
 neighborhood in the exhaustion and all sufficiently large $k$;
 \item $z_k=F_{s_a,s_k}(z)$, uniformly for
 $\abs z\le e^{-\gamma s_a}$, for any sufficiently large $a$ and
 all $k\ge a$.
\end{enumerate}
In the first case the threshold may depend on the compact; in the
second it is uniform in $a$.  Continuation and composition of the
transition germs give $z_{k+1}=F_k(z_k)$ in both cases.
The chart estimates and \eqref{eq:unif:decay-parameters} imply,
eventually,
\begin{equation}
 \abs{z_k}\le e^{-5\nu s_k},\qquad
 \zeta_k:=H_k(z_k),\quad
 \abs{\zeta_k}\le e^{-2\nu s_k}.
 \label{eq:unif:small-orbits}
\end{equation}
In regime (ii), the case $k=a$ uses $F_{s_a,s_a}=\Id$ and
$\gamma>10\nu$.

Choose $h$ so small that, with $\delta=e^h-1$,
\begin{equation}
 \delta<\tfrac1{10},\qquad
 A_d\delta<1,\qquad C_d\delta<\nu/4,
 \label{eq:unif:macro-choice}
\end{equation}
enlarging $A_d,C_d$ first to cover the finitely many fixed-degree
coefficient estimates used here.  By
\eqref{eq:unif:jet-order}--\eqref{eq:unif:small-orbits}, after the
late start the defect satisfies
\begin{equation}
 r_k:=\zeta_{k+1}-G_k(\zeta_k),\qquad
 \abs{r_k}\le e^{-10s_k}.
 \label{eq:unif:orbit-defect}
\end{equation}
All constants dependent on $h,N,M$ multiply $o(s_k)$ in this argument
and are absorbed only by that late start.  In particular,
\eqref{eq:unif:macro-choice} does not require choosing $h$ using a
constant which is itself determined by $M$.

The inverse maps have the local Lipschitz bound
\begin{equation}
 \Lip\left(G_k^{-1}\bigm|_{B(e^{-\nu s_{k+1}})}\right)
 \le e^{\Delta s_k}
 \label{eq:unif:inverse-tube}
\end{equation}
at all sufficiently large indices.  To check it, the derivative of
$G_k^{-1}$ on this ball is bounded by the sum of
\[
 s_k^C e^{\Delta s_k/2}
 \quad\hbox{and}\quad
 \exp\bigl(C_d\Delta s_k+o(s_k)-\nu s_{k+1}\bigr).
\]
The first term is at most $e^{3\Delta s_k/4}$ eventually, since
$\Delta s_k\ge c_h\sqrt{s_k}$.  The second tends to zero by
\eqref{eq:unif:macro-choice}; their sum is at most
$e^{\Delta s_k}$ after increasing the starting index.  Integrating
the derivative along segments in this convex ball gives
\eqref{eq:unif:inverse-tube}.

\subsection{Backward approximation inside valid tubes}

\begin{lemma}[Convergence of corrected coordinates]
\label{lem:unif:backward-limit}
For either raw sequence above and every sufficiently large $k$,
\[
 \xi_k=\lim_{l\to\infty}G_{k,l}^{-1}\zeta_l
\]
exists locally uniformly in the indicated variables.  It is
holomorphic and satisfies
\begin{equation}
 \abs{\xi_k-\zeta_k}\le e^{-6s_k},\qquad
 \xi_{k+1}=G_k\xi_k.
 \label{eq:unif:shadow}
\end{equation}
All applications of \eqref{eq:unif:inverse-tube} take place inside
its stated balls.
\end{lemma}

\begin{proof}
For a terminal index $l$ put $y_l^{(l)}=\zeta_l$ and
$y_i^{(l)}=G_i^{-1}y_{i+1}^{(l)}$ for $i<l$.
Simultaneously establish, by backward induction, that the comparison
segments at step $i$ lie in $B(e^{-\nu s_{i+1}})$ and that
\begin{equation}
 \abs{y_k^{(l)}-\zeta_k}
 \le \sum_{i=k}^{l-1}
       e^{-10s_i}\exp(s_{i+1}-s_k).
 \label{eq:unif:finite-error-sum}
\end{equation}
The terminal assertion is immediate.  At a backward step compare
$G_i^{-1}y_{i+1}^{(l)}$ with
$G_i^{-1}G_i\zeta_i=\zeta_i$.
Both arguments are close to $\zeta_{i+1}$: the first by the inductive
bound and the second by \eqref{eq:unif:orbit-defect}.  Since
$\abs{\zeta_{i+1}}\le e^{-2\nu s_{i+1}}$, they and their joining
segment lie in the required larger tube whenever the error sum has
the bound asserted below.  Applying
\eqref{eq:unif:inverse-tube} then gives the next error sum.

Indeed $s_{i+1}\le(1+\delta)s_i$ and the grid eventually has spacing
at least one, so
\[
 \sum_{i\ge k}e^{-10s_i+s_{i+1}-s_k}
 \le e^{-s_k}\sum_{i\ge k}e^{-(9-\delta)s_i}
 \le C e^{-(10-\delta)s_k}\le e^{-6s_k}.
\]
As $\nu<1/10$, this is much smaller than the difference between the
two tube radii.  The estimates therefore close the simultaneous
induction without evaluating an inverse outside its proved domain.

For completeness, consecutive terminal approximations are Cauchy
by the same reasoning.  At level $l$ their difference is at most
$e^{\Delta s_l}e^{-10s_l}$, and propagation down to $k$ bounds it by
$e^{-10s_l+s_{l+1}-s_k}$.  These bounds are summable in $l$.
They are uniform on the compact sets in either regime, so the limits
are holomorphic.  The finite identities give
$\xi_{k+1}=G_k\xi_k$, and the infinite error sum proves
\eqref{eq:unif:shadow}.
\end{proof}

\subsection{A coherent injective map on the manifold}

\begin{proposition}[Global coordinates]\label{prop:global-coordinates}
The locally uniform limit on the exhaustion
\begin{equation}
 \Psi(x)=\lim_{l\to\infty}G_{0,l}^{-1}H_lf_{s_l}(x)
 \label{eq:unif:global-map}
\end{equation}
defines an injective holomorphic map $M\to\C^n$ with $\Psi(o)=0$.
Its image $\Omega$ is open, and $\Psi:M\to\Omega$ is a
biholomorphism.  Moreover, on every compact $K\subset\Omega$,
\begin{equation}
 \sup_K\abs{G_{0,k}}\le e^{-(\sigma/2)s_k}
 \quad\hbox{eventually},
 \label{eq:unif:forward-smallness}
\end{equation}
and, at every sufficiently large $k$,
\begin{equation}
 G_{0,k}^{-1}\bigl(B(e^{-\gamma's_k})\bigr)\subset\Omega.
 \label{eq:unif:inverse-ball-inclusion}
\end{equation}
\end{proposition}

\begin{proof}
On a fixed exhaustion compact choose an index $k$ after which regime
(i) is available.  Lemma~\ref{lem:unif:backward-limit} gives $\xi_k$
there, and composing it with the single fixed polynomial
automorphism $G_{0,k}^{-1}$ gives the limit
\eqref{eq:unif:global-map}.  On overlaps the inverse charts agree by
continuation, and \eqref{eq:unif:shadow} makes the limits agree.
Thus these maps define a holomorphic $\Psi$ on all of $M$.  All maps
fix zero in charts, so $\Psi(o)=0$.

Suppose $\Psi(x)=\Psi(y)$ for two fixed points.  For every large $k$,
$G_{0,k}\Psi=\xi_k$ on a compact containing them, and hence
\eqref{eq:unif:shadow} implies
\[
 \abs{H_kf_{s_k}(x)-H_kf_{s_k}(y)}\le2e^{-6s_k}.
\]
On the tiny raw-coordinate ball containing these points, write
$H_k(z)=A_kz+E_k(z)$, where $A_k=H_k'(0)$ has conorm at least
$s_k^{-C}$ and $\norm{DE_k}$ is bounded by $\exp(o(s_k))$ times
the radius.  Integrating only $DE_k$ along the segment and retaining
the fixed linear term $A_k$ gives
\[
 |H_k(z)-H_k(w)|\ge
 \bigl(s_k^{-C}-\sup\norm{DE_k}\bigr)|z-w|.
\]
The nonlinear error is eventually less than $s_k^{-C}/2$, so
\begin{equation}
 \abs{f_{s_k}(x)-f_{s_k}(y)}\le e^{-5s_k}
 \label{eq:unif:raw-separation}
\end{equation}
eventually.

If $x\ne y$, choose a small relatively compact initial metric ball
about $x$ whose closure excludes $y$.  Map the straight segment
between the two raw coordinates by $\Phi_{s_k}$, obtaining a path
from $x$ to $y$.  The segment stays in the uniform chart ball.
Up to its first exit from the fixed ball about $x$, the compact
metric comparison $e^{-2s_k}g\le h_{t(s_k)}$ from
Proposition~\ref{prop:ricci-windows}, and the uniform chart upper
bound $\Phi_{s_k}^*h_{t(s_k)}\le Cg_{\mathrm{Eucl}}$, bound its
initial-metric length by
$Ce^{s_k}e^{-5s_k}$.  This tends to zero and cannot reach the
boundary of that fixed ball.  The contradiction proves injectivity.
This argument uses immersed charts only on their indicated sheet;
it requires no global injectivity of $\Phi_{s_k}$.

Nonsingularity follows directly from the same quantitative estimates.
Fix a point and two coordinate neighborhoods $U\Subset V$ whose
closures lie in one compact exhaustion set.  The estimates of
regime (i) hold uniformly on $\overline V$ for every sufficiently
large $k$.  Differentiating $\Phi_{s_k}f_{s_k}=\Id$, and using
\eqref{chart:metric-decay} and
\eqref{chart:uniform-chart-metric}, gives
\[
 |df_{s_k}(v)|\ge c e^{-s_k}|v|_g
 \quad\text{on }\overline V.
\]
The coefficient bounds on $H_k$, its polynomially bounded inverse
linear part, and $|f_{s_k}|\le e^{-\sigma s_k}$ therefore imply
\[
 |d(H_kf_{s_k})(v)|\ge c s_k^{-C}e^{-s_k}|v|_g.
\]
On the other hand, the shadow bound
$|\xi_k-H_kf_{s_k}|\le e^{-6s_k}$ on $\overline V$ and Cauchy's
estimate on the fixed pair $U\Subset V$ bound its differential on
$U$ by $C e^{-6s_k}|v|_g$.  Thus
\[
 |d\xi_k(v)|\ge
 \bigl(c s_k^{-C}e^{-s_k}-C e^{-6s_k}\bigr)|v|_g>0
 \qquad(v\ne0)
\]
for all sufficiently large $k$.  Since $\xi_k=G_{0,k}\Psi$ and
$G_{0,k}$ is a polynomial automorphism, $d\Psi$ is nonsingular.
The holomorphic inverse theorem now makes the image $\Omega$ open;
the local inverses agree by the proved injectivity, giving a
holomorphic inverse on $\Omega$.  Consequently
$K'=\Psi^{-1}(K)$ is compact whenever $K\subset\Omega$ is compact.  On $K'$, the common chart bound
$\abs{f_{s_k}}\le e^{-\sigma s_k}$, the coefficient bound on $H_k$,
and \eqref{eq:unif:shadow} give
\[
 \sup_K\abs{G_{0,k}}
 \le \exp\bigl(-\sigma s_k+o(s_k)\bigr)+e^{-6s_k},
\]
which proves \eqref{eq:unif:forward-smallness}.

It remains to prove the ball inclusion.  Fix a sufficiently large
$k$ and use regime (ii) with $a=k$.  For large $l$, the inverse chart
$f_{s_l}$ is defined on the compact image under $\Phi_{s_k}$ of the
closed ball $\overline B(e^{-\gamma s_k})$.  The identity
$f_{s_l}\Phi_{s_k}=F_{s_k,s_l}$ holds first as a germ at zero and
then throughout that ball by holomorphic continuation.  Therefore
Lemma~\ref{lem:unif:backward-limit} gives
\begin{equation}
 \sup_{\abs z\le e^{-\gamma s_k}}
 \abs{G_{0,k}\Psi\Phi_{s_k}(z)-H_k(z)}
 \le e^{-6s_k}.
 \label{eq:unif:chart-shadow}
\end{equation}
Write $A=H_k'(0)$, $r=e^{-\gamma s_k}$, and
\[
 T_k=G_{0,k}\Psi\Phi_{s_k},\qquad E_k(z)=T_k(z)-Az.
\]
The map $T_k$ is holomorphic on the fixed chart ball, and fixes zero.
Its shadow estimate and the coefficient bound on $H_k$ give
\[
 E_k(0)=0,\qquad \norm{A^{-1}}\le s_k^C,\qquad
 \sup_{\abs z\le r}|E_k(z)|
 \le \exp(o(s_k))r^2+e^{-6s_k}.
\]
Cauchy's estimate on the smaller ball then gives
\[
 \sup_{\abs z\le r/2}\norm{A^{-1}DE_k(z)}
 \le C_n s_k^C\bigl(\exp(o(s_k))r+e^{-6s_k}/r\bigr)
 \longrightarrow0.
\]
For large $k$ this supremum is at most $1/4$.  Moreover,
$\gamma'>\gamma$ implies
$|A^{-1}w|\le r/4$ whenever $|w|\le e^{-\gamma's_k}$,
after increasing the starting index once more.  For such a target
$w$, the map
\[
 z\longmapsto A^{-1}w-A^{-1}E_k(z)
\]
is $1/4$-Lipschitz on $\overline B(r/2)$ and maps this closed ball
into $B(3r/8)$, because $E_k(0)=0$.  Its successive iterates are
Cauchy, stay in the closed ball, and converge to a fixed point $z$.
The fixed-point identity is $T_k(z)=w$.  Therefore
$B(e^{-\gamma's_k})\subset G_{0,k}(\Omega)$, which is precisely
\eqref{eq:unif:inverse-ball-inclusion}.
\end{proof}

\subsection{The image is all of affine space}
\label{sec:surjectivity}

We first record the polynomial growth estimate used in the final step.
For a vector polynomial $P$ of degree at most $m$ and $0<R_1<R_2$,
\begin{equation}
 \sup_{B(R_2)}\abs P
 \le (R_2/R_1)^m\sup_{B(R_1)}\abs P.
 \label{eq:unif:polynomial-growth}
\end{equation}
To verify this, restrict each scalar dual pairing of $P$ to a
complex line through zero.  The maximum principle in the exterior
of the radius-$R_1$ disc, applied to the polynomial divided by
$z^m$ including its removable value at infinity, gives the scalar
bound.  Taking the supremum over unit dual vectors and lines proves
\eqref{eq:unif:polynomial-growth}.

\begin{proof}[Proof of Theorem~\ref{thm:uniformization}]
The preceding analytic and chart constructions, followed by
Propositions~\ref{prop:polynomial-models} and
\ref{prop:global-coordinates}, give a biholomorphism $\Psi:M\to\Omega$
onto an open set containing zero, with the forward smallness and
inverse-ball inclusion proved above.  It remains to show that
$\Omega=\C^n$.

Suppose
$R_*=\sup\{R>0:B(R)\subset\Omega\}<\infty$.
It is positive because $0\in\Omega$ and $\Omega$ is open.
By \eqref{eq:dyn:forward-degree} there is an unbounded subsequence
with $\deg G_{0,k}\le Ds_k$ for some finite $D$.
Choose
$0<R_1<R_*<R_2$ so close to $R_*$ that
\begin{equation}
 D\log(R_2/R_1)<\sigma/2-\gamma'.
 \label{eq:unif:radii-choice}
\end{equation}
The closed smaller ball lies compactly in $\Omega$: it lies inside
some still larger centered ball of radius below $R_*$.
By \eqref{eq:unif:forward-smallness},
\eqref{eq:unif:polynomial-growth}, and
\eqref{eq:unif:radii-choice}, the chosen subsequence satisfies
\[
 \sup_{B(R_2)}\abs{G_{0,k}}
 \le \exp\left[
   -\frac{\sigma}{2}s_k+Ds_k\log(R_2/R_1)\right]
 < e^{-\gamma's_k}
\]
for all sufficiently large indices.
Equation~\eqref{eq:unif:inverse-ball-inclusion} then implies
$B(R_2)\subset\Omega$, contradicting the definition of $R_*$.
Therefore $R_*=\infty$, and $\Omega=\C^n$.

The map $\Psi$ in \eqref{eq:unif:global-map} is consequently a
biholomorphism $M\to\C^n$.
\end{proof}

The surjectivity step used the degrees
of the forward compositions $G_{0,k}$ on their available
subsequence; no degree estimate for their inverses, or uniform
instantaneous contraction ratio, is required.
